\chapter{Introduzione e scopo del documento}

\section{Il problema}

Il \emph{Non-Intrusive Load Monitoring} (NILM) consiste nello stimare quali elettrodomestici
sono accesi in ogni istante, disponendo del solo consumo \textbf{aggregato} dell'abitazione
misurato da un unico dispositivo IoT. Non esistono etichette temporali: nessuno ha registrato
quando la lavatrice è stata effettivamente avviata, quindi il problema è per costruzione
\textbf{non supervisionato} e non è verificabile rispetto a una verità di riferimento.

Il modello di misura è additivo:
\begin{equation}
  y(t) \;=\; \sum_{i=1}^{N} P_i \, x_i(t) \;+\; \eta(t),
  \label{eq:modello-misura}
\end{equation}
dove $y(t)$ è la potenza aggregata letta dal contatore, $P_i$ la potenza assorbita
dall'elettrodomestico $i$-esimo quando è in funzione, $x_i(t) \in \{0,1\}$ indica se è acceso e
$\eta(t)$ raccoglie rumore di misura e carichi non modellati.

Invertire la~\eqref{eq:modello-misura} è un problema mal posto: molte combinazioni diverse di
elettrodomestici producono la stessa somma. Con dieci apparecchi ci sono $2^{10}$ configurazioni
possibili per ogni istante, e un plateau da $1500\unit{W}$ può essere spiegato da una
lavastoviglie, da un climatizzatore o dalla somma di due carichi minori in modo perfettamente
equivalente rispetto al solo segnale.

\section{Il ruolo del questionario}

L'informazione che rende il problema trattabile non viene dal segnale ma dal
\textbf{questionario} compilato dall'utente, raccolto attraverso una webapp dedicata descritta nel
capitolo~\ref{ch:questionario}. Per ogni elettrodomestico presente in casa si dichiara:

\begin{itemize}[nosep,leftmargin=1.4em]
  \item quanta potenza assorbe quando è in funzione (tipica, minima e di picco);
  \item quante volte alla settimana viene usato (minimo e massimo);
  \item quanto dura tipicamente un ciclo (minimo e massimo);
  \item quante ore al giorno resta acceso;
  \item in quale fascia oraria viene tipicamente avviato o usato;
  \item in quali mesi dell'anno viene usato, cioè se è stagionale;
  \item se è invece un apparecchio sempre alimentato, come frigoriferi e congelatori;
  \item quanto l'utente è sicuro di ciò che ha dichiarato.
\end{itemize}

A queste si aggiungono due dati della casa nel suo complesso: la potenza contrattuale e il carico
base sempre presente.

Tutti i modelli descritti in questo documento traducono queste dichiarazioni in
\textbf{penalità nella funzione obiettivo}, non in vincoli rigidi. La scelta è deliberata e
attraversa l'intero progetto: il questionario è una fonte di informazione utile ma
approssimativa, compilata a memoria, e va trattato come conoscenza a priori che il segnale può
smentire quando l'evidenza è sufficientemente forte. Un vincolo rigido, al contrario, rende un
errore di compilazione irrecuperabile.

\begin{notebox}{Principio guida}
Il questionario è un \emph{prior}, non un \emph{veto}. Ogni dichiarazione dell'utente diventa un
costo da pagare per violarla, mai un divieto assoluto. La misura di quel costo è il contenuto
principale di questo documento.
\end{notebox}

\section{Cosa contiene questo documento}

Il capitolo~\ref{ch:questionario} descrive la \textbf{webapp di raccolta dati}: come è strutturata
l'intervista al cliente, quali grandezze vengono chieste e in che forma arrivano al solver.

Il resto del documento descrive il modello attualmente in uso, \code{highs\_survey\_prior}: le variabili di
decisione, la funzione obiettivo e l'insieme completo dei vincoli. Il modello traduce ogni
dichiarazione del questionario in una penalità dell'obiettivo, secondo il principio guida enunciato
sopra.

Il capitolo~\ref{ch:obiettivo} è dedicato interamente all'\textbf{anatomia della funzione
obiettivo}: cosa rappresenta ciascun termine, perché è scalato in quel modo, quale significato
concreto ha il valore numerico del suo coefficiente e come si comporta il modello al variare di
quel valore. È la parte da leggere per capire come tarare il sistema.

\begin{notebox}{Nota sulla denominazione}
La chiave dell'approccio nel codice è \code{highs\_survey\_prior}. Il modello è costruito e risolto
invocando direttamente HiGHS attraverso la libreria \code{highspy}, un solver di programmazione
lineare mista intera libero e senza licenza.
\end{notebox}
